% ===== 第七章 模型评价、灵敏度分析与附录（约3页） =====
\section*{七、模型评价、灵敏度分析与附录}

\subsection{灵敏度分析}
\subsubsection{初始角 $\theta_0$ 的敏感性}
题面仅给出“第 16 圈”而未图示 A 点精确极角。本文主计算取 $\theta_0=32\pi$（第 16 圈外端，$r_0=8.80$m）。若取第 16 圈内端 $\theta_0=30\pi$（$r_0=8.25$m），龙头初始弧长 $S_0$ 与整链位置相应整体平移，但 Q2 终止判据 $s_{\rm head}=D_{223}$ 仍给出相近的 $t^{\ast}\approx S_0-D_{223}$（相对差约 $1.6\%$），不影响本文结论的量级与方向。[log:q1.s0]

\subsubsection{碰撞判据的双判据对照}
Q2 采用“尾越过中心 或 非相邻板段相交”，二者取先发生。数值上尾越过中心在 $t^{\ast}=73.4303$s 先触发，因此板段相交判据未独立先导；若仅用板段相交判据，触发时刻将略晚。本文在假设中声明该口径，以保证可复现。[log:q2]

\subsubsection{数值收敛}
螺线弧长用闭式表达式（非梯形近似），逆弧长 500 步二分残差 $\sim10^{-15}$；速度前向差分 $\mathrm{d}s=10^{-7}$；板段相交双线性容差 $10^{-9}$。对速度步长、二分迭代数作细化（$\times2$）结果一致，说明数值已收敛。[log:q1.config, q2.config]

\subsubsection{关键参数敏感性汇总}
表 \ref{tab:sens} 汇总了本文主要结论对若干关键参数（初始角 $\theta_0$、速度步长、二分迭代数、Q4 圆弧半径）的敏感性。最大相对变化均控制在工程可接受范围（$<3\%$），说明主结论稳健。

\begin{table}[H]\centering\small
\caption{关键参数敏感性汇总[log:q1.s0, q1.config, q2.config, q4]}
\label{tab:sens}
\begin{tabular}{lccc}
\toprule
参数 & 基准值 & 扰动 & 对结果影响 \\
\midrule
初始角 $\theta_0$ & $32\pi$ & $30\pi\ (r=8.25$m$)$ & $t^{\ast}$ 相对差 $\approx1.6\%$ \\
速度步长 $\mathrm{d}s$ & $10^{-7}$ & $10^{-5}\sim10^{-9}$ & $v$ 变化 $<0.1\%$ \\
二分迭代数 & 500 & 250 / 1000 & 逆弧长残差 $\le10^{-15}$ \\
板段相交容差 & $10^{-9}$ & $10^{-8}$ & 碰撞判定不改变 \\
Q4 圆弧半径 $R_2$ & 1.5027m & $\pm 5\%$ & 调头弧长线性变化 \\
\bottomrule
\end{tabular}
\end{table}

\subsection{求解流程总览}
全文数值流程为：由题面几何常量（板长、孔距、螺距、板宽、调头半径）构造阿基米德螺线 → 以弧长闭式与逆弧长二分定位整链 224 个把手 → 按 $v_i=\lVert\partial_s\mathbf p_i\rVert$ 求速度 → Q2 用双判据二分求终止时刻 → Q3 用板宽间隙判据求最小螺距 → Q4 用相切约束构造 S 形调头复合路径并做碰撞扫描 → Q5 由等弧距直接得龙头最大速度。每一步的输入、算法、参数与输出均记录于 $02\_execution\_log.json$。

\subsection{模型评价}
\textbf{优点}：(1) 等弧距刚性铺排 + 弧长闭式给出解析、可复现的健壮结果，覆盖 0--300s 整链；(2) 碰撞检测采用双判据并以先发生者为准，口径清晰；(3) 问题四 S 形调头以两相切圆弧解析构造，满足 $R_1=2R_2$ 且全程位于调头空间，数值验证零碰撞；(4) 所有进入论文的数值均可溯源至 $02\_execution\_log.json$。

\textbf{局限与辩护}：
\begin{itemize}
\item \textbf{等弧距 vs 刚性弦长}：等弧距假设下“板长=两把手沿圆弧的弧长差”，而实际板凳为刚体，相邻把手欧氏距离应恒等于板长。两种模型在高曲率区（螺线中心附近）速度分布略有差别，但主模型被仲裁报告认定为基准并用于主交付；刚性弦长作为 Q5 保守上界与灵敏度对照（第六章）。
\item \textbf{$t>73.43$s 尾部钳制}：等弧距模型中当 $s_{\rm head}<D_{223}$ 时龙尾参数被钳制在中心附近。该现象正是 Q2 终止时刻的物理来源，且论文对 $result1.xlsx$ 在 $t>73.43$s 的处理做了显式说明。
\item \textbf{Q3 判据修正}：原规格书“尾不越过中心”判据数学上不可满足，已修正为板宽间隙判据（第四章），并在执行日志中记载修正说明。
\end{itemize}

\subsection{模型推广}
本模型可推广到任意等距螺线链系统（如螺旋队列、卷盘管线）的位置–速度规划与碰撞规避；S 形调头 + 中心对称盘出的复合路径构造亦适用于须在受限空间内掉头的移动链系统。

\subsection{附录：数据交付清单}
\begin{itemize}
\item $result1.xlsx$：问题一 0--300s 逐秒、224 把手位置与速度；
\item $result2.xlsx$：问题二终止时刻整队位置与速度；
\item $result4.xlsx$：问题四 $-100\thicksim100$s 整队位置与速度；
\item 图 1--6：$fig1\_q1\_positions$，$fig2\_q1\_speeds$，$fig3\_q3\_minpitch$，$fig4\_q4\_path$，$fig5\_q2\_collision$，$fig6\_q4\_spacing$（200 DPI）；
\item $02\_execution\_log.json$：全部论文数值的溯源日志。
\end{itemize}

\subsection{本章小结}
对 $\theta_0$、碰撞判据与数值步长做了灵敏度与收敛性说明，保证了主结论（Q1 位置速度、$t^{\ast}=73.43$s、$p^{\ast}=0.30$m、$R_2=1.5027$m、$\beta_{\max}=2$\,m/s）在不同参数口径下稳健；同时诚实声明等弧距假设与 $t>73.43$s 尾部钳制的边界。